针对啤酒化合物到感官画像的预测、关键因素识别与风格迁移三个问题，本文先重建评分端点，再分别建立预测、解释与决策模型。
\par
附件 4 的感官值是逐人中心化后的连续值，不是 0--6 原始分。本文据公开的 2197 条 0--6 评分（16 名品评员，每款酒 3--14 人）重建了 50 个目标的面板均分与七级分布；对8.xlsx的70个参考字段逐项核查：37个符合对数规则、1个恒等对应，另32个字段含80处差异。附件中显式缺失仅 73 格，另有 18.7\% 的输入单元格为负值，其含义需按字段变换规则判断，二者必须分开处理。
\par
第一问按酒三折五次验证。极端随机树与梯度提升树的等权融合在核心 30 项上的面板均分平均绝对误差为 0.50222 分，明显优于训练均值（0.60450 分），却只比风格中位数基线（0.50969 分）低 0.00747 分，差值的 95\% 重采样区间 $[-0.01914,0.00418]$ 含 0：固定预测比较不足以确认融合模型优于已知风格基线；补充的同参数ET对照中，风格加化学相对仅风格的核心MAE低0.01426分，但相对风格中位数的区间仍跨0。七级分布模型的累计分布平方距离为 0.02178（总体基线 0.02820），面板众数一致率 80.52\%，但非零众数一致率仅 24.33\%，总体一致率主要由 78.03\% 的零等级贡献。保形区间逐属性平均覆盖 90.06\%、宽 1.7342 分，按品评人数分组后 3--4 人组降至 74.00\%。缺失机制实验显示，系统性低值缺失使核心误差增加 0.0273 分，约为随机缺失 20\%（0.0139 分）的两倍，不同缺失机制和缺失比例需要共同解释。
\par
第二问以 $|r|\ge0.7$ 构建 231 节点、109 条边的共表达图。图正则模型在单折平均口径下略差于无图（0.537035 对 0.536911），在聚合预测口径下略好（0.524865 对 0.526148），最终模型取 $\lambda=0$；网络主要用于描述共变结构；150 条 Top-3 候选中 36 条满足“频率 $\ge0.5$、至少 5 折入选、置换贡献为正”的条件，将 80 个异常单元格置缺后 Top-3 保留 97.64\%，仅用一致字段与有单位登记字段的32项交集时降至36.18\%。
\par
第三问对 Tripel、Kriek、Lambic 各生成 40 条谱。受约束局部凸组合把 Tripel 的平均目标风格概率提高到 0.3358（重采样 0.2992），判为目标风格的比例为 0.695（重采样 0.565），但分布距离反而更大（0.1158 对 0.1034），说明“更像目标风格”与“更接近真实分布”不是一回事。迁移在 55 条记录上搜索可行 51 条、独立代理支持 11 条；严格浓度协议下最大 Tripel 份额全为 0，且 55 条记录均找到正余量不可行证书，阻塞主要来自 DMS 与反式-2-壬烯醛；改用阈值超额协议后，参考阈值下30条记录允许正份额，但没有候选通过联合代理复核；2倍阈值敏感性下34条允许正份额，其中9条重复划分记录获得代理支持，并未经过实际制备与盲评。
\par
结果显示：风格基线与复杂模型的误差接近；关键因素应按稳定支持与探索性候选分层解释；阈值约束与代理判断显著影响迁移候选，实际效果需要制备和盲评。
